
\textbf{Session 2, February 16: Grand theories of war and peace}

\begin{itemize}
    \item [***] \parencite{buenodemesquitaPrinciplesInternationalPolitics2013}[Chapter 5]
    \item \parencite{levyCAUSESWARCONDITIONS}
    \item \parencite{lakeWhyIsmsAre2011}
    \item \parencite{lakeChallengesLiberalOrder2021}
    \item \parencite{milnerMythEclecticIR2023}
    \item \parencite{OneWorldRival}
\end{itemize}

\newpage

\section{\cite{buenodemesquitaPrinciplesInternationalPolitics2013} Chapter 5}

\begin{itemize}
    \item War is generally inefficient compared to reaching a negotiated settlement.
    \item The risks of war may be exacerbated by uncertainty, making credible commitments, or winner-takes-all situation.
    \item Some important neorealist hypotheses about international peace and stability do not
    flow logically from that theory’s core assumptions. Other important neorealist hypotheses are falsified by the empirical record.
    \item The link between power transition theory’s assumptions and hypotheses is more clear-cut than is true in the case of neorealism.
    \item Some core power transition hypotheses follow from the theory’s logic and are supported by evidence; others either do not follow logically or are not borne out by evidence.
    \item The power transition theory outperforms neorealism when it comes to offering a structural explanation for war, but its weaknesses point to a need to reconsider its approach to domestic politics.
\end{itemize}

Consider that estimates of war-related deaths per year in the past 100 years vary approximately from 1 to 2 million.

\subsection{War and Uncertainty}
The ex ante problems that seem to lead to war can be reduced to three main factors: (1) uncertainty, (2) commitment problems, and (3) indivisibility of issues (Fearon 1995; Slantchev 2003).

\begin{itemize}
    \item [1] Absence of \textbf{credible commitments} between the sides.
    \item [1.1] The inability to \textbf{commit to negotiating} patiently rather than seizing the first-strike advantage
    exacerbates the risk of war (Powell 1999; Slantchev 2003)
    \item [1.2] The idea is that the Israelis give up
    some land in exchange for a guarantee of peace. An alternative formulation reverses the sequence of concessions—the Palestinian Authority disarms its factions that engage in violence against Israel in return for land concessions. Either approach is likely to fail because each induces a fundamental commitment problem.
    \item [] \textbf{Land for peace} suffers from a time inconsistency problem.
    \item [2] Dispute is about an \textbf{indivisible objective} --> No settlement can produce a compromise in this case of indivisibility, and so there is a severe risk of war. If control over all the land currently held by Israel
    is truly its irreducible objective, then the land is viewed as indivisible by Hamas.
    \item [3] This \textbf{uncertainty about the true probability of victory} or defeat might arise out
    of private (secret) information about the morale of their fighting forces, about knowledge of
    outside help, about the quality of their war plans, or a host of other possible factors.
\end{itemize}

\subsection{Realist Theory}
Neorealist theorists believe that the distribution of power in the international system is a
major factor in determining whether international affairs are stable or unstable. Stability
refers to circumstances in which the sovereignty of key states is preserved (Gulick 1955).
Instability refers to changes in the composition of the international system, especially
changes involving the disappearance or emergence of “key” states following large wars. Key
states are those whose assistance might be necessary to counteract a threat from a rival coali-
tion of states.

In addressing war and instability, neorealist theories start with the following
four assumptions:

\begin{enumerate}
    \item International politics is anarchic. $\Rightarrow$ there is no supranational authority that can enforce agreements between states.
    \item States, as rational unitary entities, are the central actors in international politics. $\Rightarrow$ domestic politics does not influence international politics.
    \item States seek to maximize their security above all else and consider other factors only after security is assured. $\Rightarrow$ states are not willing to trade away their security for other benefits. Other possible goals, such as wealth, are pursued only after security has been assured (Waltz 1979, 126).
    \item States seek to increase their power as long as doing so does not place their security at risk. $\Rightarrow$ states accept being weaker than they might otherwise be if the pursuit of greater power would place their security at risk.
\end{enumerate}

Most important neorealist conclusions
\begin{enumerate}
    \item Bipolar systems are more stable than multipolar systems.
    \item States engage in balancing behavior so that power becomes more or less equally divided
    among states over time.
    \item States mimic, or echo, one another’s behavior.
\end{enumerate}

\newpage

\section{How Well Does Neorealism Do in Explaining War and Instability?}

\subsection{Bipolarity and Stability}
The argument that bipolar structures are more stable than multipolar ones is built on the claim that there is more uncertainty in a multipolar system than in a bipolar one.

Bipolarity is thought to alleviate two of the conditions we have discussed that are thought to increase the risk of war: uncertainty and commitment problems.

There are several problems with the argument that because bipolar systems encompass less
uncertainty than multipolar systems they yield greater stability. This argument is not implied logically by the four key assumptions of neorealism.

\begin{enumerate}
    \item First, variation in the response to uncertainty violates the unitary rational actor assumption, which contends that important choices in international politics are driven by structural factors, not by considerations internal to the state.
    \item Second, if different decision makers respond to uncertainty in different ways, then there is no reason to expect any empirical relationship between bipolarity and stability at all.
    \item Third, if leaders differ in how they respond to uncertainty, then the third hypothesis (that states mimic one another) does not hold up.
    \item taking the neorealist assumptions into account, we can show that it is logically true that more distributions of power are stable in a multipolar world than in a bipolar world. 
\end{enumerate}

\subsection{Bipolarity and Stability 2}
Stability in a multipolar world is not limited to the situation of exact power equality. This is in marked contrast to the bipolar world, in which stability can be achieved only through a perfectly equal division of resources between the two blocs.

Thus, it appears that the first neorealist hypothesis, that bipolarity promotes stability and multipolarity promotes instability, is logically flawed given the assumptions of the theory.

Further, it seems that a true balance of power is essential for stability in a bipolar world, but not in a multipolar one, contradicting the second hypothesis. A vast array of power distributions can produce stability. In a multipolar environment, an exact balance of power is irrelevant.

Either the balance of power does not matter in multipolarity or the term is defined to mean any system in which each actor is essential. In the latter case, so many systems qualify as a balance-of-power system that the concept doesn’t have much meaning in our search for answers.

\subsection{History and Neorealist Empirical Claims}

\begin{enumerate}
    \item The international system was multipolar in structure from 1648 until the defeat of Germany in 1945. Thus, between 1648 and 1945 there were many major powers with no one or two predominating. $\Rightarrow$ The bipolar system began in 1945 and lasted until about 1989.
    \item Using Jack Levy’s (1983) classification of great powers since 1492, we can determine whether the longevity of the bipolar system was comparatively long or short. $\Rightarrow$
    \item Another way to evaluate the stability of the international system is to examine the frequency of wars among the most influential states in the world
    \item The UN Security Council (UNSC), then, by institutionalizing a multipolar decision-making structure, has provided a counterweight to bipolarity. This too may be an explanation for the long peace.
    \item The average interval between general wars is thirty-four years; the longest interval is ninety-nine years, much longer than the so-called long peace. It appears, then, that the bipolar, long peace cannot be considered unusually long after all and that bipolarity is not a superior treatment for the disease of major power war than is multipolarity.
\end{enumerate}

\subsection{Historical Record of Essential and Inessential States}

An essential state is any state that can join a losing coalition and, by dint of its membership and the power it can bring to bear, turn that coalition into a blocking coalition or a winning coalition.

Inessential states are states that are unable to turn even one losing alliance into a winning or blocking coalition.

\begin{enumerate}
    \item Essential states never become inessential.
    \item Essential states are never eliminated from the international system.
    \item Inessential states never become essential states.
    \item Inessential states are always eliminated from the international system.
\end{enumerate}

Each of these propositions is historically false.

\subsection{Balance of Power and Neorealism}
Perhaps the most common claim among neorealist theorists is that war occurs only when both sides believe that their chance of winning is greater than 50 percent. This is a partial statement of what has come to be known as the balance-of-power theory. It invokes the idea that war is the product of mutual optimism by rival states, meaning that they overestimate their chances of victory

If the neorealist hypothesis were correct, we should observe no wars initiated by the weaker side. The evidence fails to support the hypothesis.

\newpage